% ============================================================
% PREAMBLE - Configuración general del PDF
% Informe Integral de Pruebas - ECommerce Sigchos
% Basado en la plantilla institucional de Manual-Usuario-Sigchos
% ============================================================

% --- PAQUETES NECESARIOS ---
\usepackage{fancyhdr}
\usepackage{graphicx}
\usepackage{xcolor}
\usepackage{tabularx}
\usepackage{multirow}
\usepackage{makecell}
\usepackage{lastpage}
\usepackage{tcolorbox}
\usepackage{fontspec}
\usepackage{setspace}
\usepackage{titlesec}
\usepackage{float}
\usepackage{caption}
\usepackage{hyperref}
\usepackage{enumitem}
\usepackage{fontawesome5}
\usepackage{etoolbox}

% --- AJUSTE DE TABLAS (longtable de Quarto/Pandoc) ---
\AtBeginEnvironment{longtable}{%
  \fontsize{9pt}{11pt}\selectfont%
  \setlength{\tabcolsep}{4pt}%
  \renewcommand{\arraystretch}{1.15}%
  \sloppy%
}

% --- Cargar librería de tcolorbox ---
\tcbuselibrary{skins,breakable}

% --- Suprimir portada automática de LaTeX ---
\renewcommand{\maketitle}{}

% --- FUENTE: Times New Roman 12pt ---
\setmainfont{Times New Roman}

% --- COLORES INSTITUCIONALES ---
\definecolor{municipioVerde}{RGB}{49,192,144}
\definecolor{municipioBlanco}{RGB}{255,255,255}
\definecolor{espeblue}{RGB}{0,51,102}
\definecolor{infobg}{RGB}{219,237,252}
\definecolor{infoborder}{RGB}{0,102,204}
\definecolor{importantbg}{RGB}{255,243,205}
\definecolor{importantborder}{RGB}{255,193,7}

% --- VARIABLES DEL DOCUMENTO ---
\newcommand{\docCode}{SIGCHOS-PRU-V1-2026-001}

% --- INTERLINEADO: 1.5 (técnico) ---
\onehalfspacing

% --- ALINEACIÓN: Izquierda (sin justificar) ---
\raggedright

% --- CAJAS DE INFORMACIÓN ---
\newtcolorbox{infobox}[1][Información]{
  colback=infobg, colframe=espeblue,
  leftrule=4pt, rightrule=0pt, toprule=0pt, bottomrule=0pt,
  left=6pt, right=6pt, top=4pt, bottom=4pt,
  breakable,
  before upper={{\textcolor{espeblue}{\faInfoCircle} \textbf{#1:}}\enspace},
  fontupper=\small
}

\newtcolorbox{importantbox}[1][Importante]{
  colback=importantbg, colframe=importantborder,
  leftrule=4pt, rightrule=0pt, toprule=0pt, bottomrule=0pt,
  left=6pt, right=6pt, top=4pt, bottom=4pt,
  breakable,
  before upper={{\textcolor{importantborder}{\faExclamationTriangle} \textbf{#1:}}\enspace},
  fontupper=\small
}

% --- NOMBRES EN ESPAÑOL (APA) ---
\captionsetup[figure]{
  name=Figura,
  labelsep=period,
  justification=raggedright,
  singlelinecheck=false,
  font={small},
  labelfont={bf}
}

\captionsetup[table]{
  name=Tabla,
  labelsep=period,
  justification=raggedright,
  singlelinecheck=false,
  font={small},
  labelfont={bf}
}

% --- FORMATO DE TÍTULOS ---
% Capítulo (nivel 1) - en negrita, con línea negra abajo
\titleformat{\chapter}[block]
  {\normalfont\Large\bfseries}
  {\thechapter.}
  {0.5em}
  {\MakeUppercase}
  [\vspace{2pt}{\color{black}\titlerule}]
\titleformat{name=\chapter,numberless}[block]
  {\normalfont\Large\bfseries}
  {}
  {0em}
  {\MakeUppercase}
  [\vspace{2pt}{\color{black}\titlerule}]
\titlespacing*{\chapter}{0pt}{-20pt}{15pt}

% Sección (nivel 2) - negrita, sin línea negra abajo
\titleformat{\section}[block]
  {\normalfont\large\bfseries}
  {\thesection.}
  {0.5em}
  {\MakeUppercase}
\titleformat{name=\section,numberless}[block]
  {\normalfont\large\bfseries}
  {}
  {0em}
  {\MakeUppercase}
\titlespacing*{\section}{0pt}{14pt}{6pt}

% Subsección (nivel 3) - negrita normal
\titleformat{\subsection}
  {\normalfont\normalsize\bfseries}
  {\thesubsection.}{0.5em}{}
\titlespacing*{\subsection}{0pt}{10pt}{4pt}

% --- ENCABEZADO: Tabla de 3 columnas ---
\pagestyle{empty}
\fancyhf{}

\fancyhead[C]{%
  \renewcommand{\arraystretch}{1.1}%
  \setlength{\tabcolsep}{3pt}%
  \begin{tabularx}{\textwidth}{|m{2.3cm}|>{\centering\arraybackslash}X|m{5.5cm}|}
    \hline
    \multirow{3}{*}{\centering\includegraphics[width=2.1cm,keepaspectratio]{imagenes/nuevo_logo.png}} &
    \multirow{3}{*}{\textbf{\fontsize{9.5pt}{11pt}\selectfont\makecell{INFORME INTEGRAL DE PRUEBAS DEL\\[1pt]SISTEMA DE COMERCIO ELECTRÓNICO}}} &
    \textbf{\fontsize{8.5pt}{10pt}\selectfont ÁREA DE TECNOLOGÍA E INNOVACIÓN} \\ \cline{3-3}
    & & \fontsize{8pt}{9pt}\selectfont Código de documento: \docCode \\ \cline{3-3}
    & & \fontsize{8.5pt}{10pt}\selectfont Página: \thepage\ de \pageref*{LastPage} \\
    \hline
  \end{tabularx}%
}

% Pie de página: código de proceso (izq) y revisión (der)
\fancyfoot[L]{\footnotesize\textit{Código de proceso: SIG.1.0}}
\fancyfoot[R]{\footnotesize\textit{Rev. ATI: 2026-nov-18}}

\setlength{\headheight}{50pt}
\setlength{\headsep}{25pt}

% Redefinir "plain" inicial para que NO TENGA ENCABEZADO en el índice
\fancypagestyle{plain}{
  \fancyhf{}
  \renewcommand{\headrulewidth}{0pt}
}

% Crear el estilo para contenido
\fancypagestyle{contentplain}{
  \fancyhf{}
  \fancyhead[C]{%
    \renewcommand{\arraystretch}{1.1}%
    \setlength{\tabcolsep}{3pt}%
    \begin{tabularx}{\textwidth}{|m{2.3cm}|>{\centering\arraybackslash}X|m{5.5cm}|}
      \hline
      \multirow{3}{*}{\centering\includegraphics[width=2.1cm,keepaspectratio]{imagenes/nuevo_logo.png}} &
      \multirow{3}{*}{\textbf{\fontsize{9.5pt}{11pt}\selectfont\makecell{INFORME INTEGRAL DE PRUEBAS DEL\\[1pt]SISTEMA DE COMERCIO ELECTRÓNICO}}} &
      \textbf{\fontsize{8.5pt}{10pt}\selectfont ÁREA DE TECNOLOGÍA E INNOVACIÓN} \\ \cline{3-3}
      & & \fontsize{8pt}{9pt}\selectfont Código de documento: \docCode \\ \cline{3-3}
      & & \fontsize{8.5pt}{10pt}\selectfont Página: \thepage\ de \pageref*{LastPage} \\
      \hline
    \end{tabularx}%
  }
  \fancyfoot[L]{\footnotesize\textit{Código de proceso: SIG.1.0}}
  \fancyfoot[R]{\footnotesize\textit{Rev. ATI: 2026-nov-18}}
}

% --- FORMATO APA: Sin sangría, espacio entre párrafos ---
\setlength{\parindent}{0pt}
\setlength{\parskip}{6pt}

% --- Listas con formato APA ---
\setlist[itemize]{leftmargin=2em, itemsep=3pt, parsep=0pt}
\setlist[enumerate]{leftmargin=2em, itemsep=3pt, parsep=0pt}
